% Generated draft: requires mathematical and source review.
\documentclass[11pt]{article}
\usepackage[margin=25mm]{geometry}
\usepackage{fontspec}
\setmainfont{Latin Modern Roman}
\usepackage{amsmath,amssymb,mathtools,unicode-math}
\setmathfont{Latin Modern Math}
\usepackage{enumitem,needspace,xurl,hyperref}
\setlength{\parindent}{0pt}
\setlength{\parskip}{5pt}
\setlength{\emergencystretch}{3em}
\newcommand{\N}{\mathbb{N}}
\newcommand{\Z}{\mathbb{Z}}
\newcommand{\Q}{\mathbb{Q}}
\newcommand{\R}{\mathbb{R}}
\newcommand{\C}{\mathbb{C}}
\newcommand{\F}{\mathbb{F}}
\newcommand{\A}{\mathbb{A}}
\newcommand{\PP}{\mathbb{P}}
\newcommand{\E}{\mathbb{E}}
\newcommand{\eps}{\varepsilon}
\newcommand{\dd}{\,\mathrm{d}}
\begin{document}
\section*{UnsolvedMath ID: 30005324}
\textbf{Source problem code:} catalogue entry
\subsection{Definitions and mathematical statement}
For a cover \(f\colon X\to Y\) belonging to the geometric enumeration referenced in the source, write \(n=\deg(f)\) for its degree. A cover is connected when its covering surface \(X\) is connected; write \(g\) for the genus of that connected surface. Denote the corresponding connected Hurwitz number provisionally by \(H^{\mathrm{conn}}_{g,n}\). This notation does not specify a counting convention: the supplied passage refers to an earlier geometric definition of \(H_{n,\ell}\) that is not included in the record.

\textbf{Requested asymptotic enumeration.} Determine asymptotic formulas for these connected Hurwitz numbers in the simultaneous limit
\[
 g\longrightarrow\infty,
 \qquad n\longrightarrow\infty.
\]
A precise asymptotic theorem must specify the underlying covers and their weights, the admissible pairs \((g,n)\), the joint growth regimes, and the meaning and uniformity of the approximation. The supplied record does not determine these choices, so it does not support a uniquely specified numerical asymptotic statement.

The source proposes, as a possible route, sufficiently precise asymptotics for the disconnected numbers \(H_{n,\ell}\) in all regimes of \(n\) and \(\ell\). Neither the definition of \(\ell\) nor its relation to genus is supplied. The passage also raises a possible first step toward a bijective proof: a combinatorial interpretation on partitions of a quantity transcribed as \texttt{C\textbackslash lambda2}. Its intended notation and definition are unresolved. The source reports a signed-bijection proof for the case \(\ell=2\) in an article by Fang, cited there as reference [3]; the record does not specify the identity proved or provide that reference's bibliographic details.
\subsection{Short English statement}
Find asymptotic formulas for connected Hurwitz numbers as both the genus of the covering surface and the degree of the cover tend to infinity. The supplied excerpt omits the original counting definition and the joint growth regimes, so these must be recovered before the request becomes a fully specified asymptotic problem. The disconnected enumeration and partition interpretation are additional source context.
\subsection{Sources}
\begin{itemize}
\item Local registry: \texttt{data/unsolvedmath/problems.json}, permanent UnsolvedMath ID \(30005324\), source problem code \texttt{catalogue entry}, title: \emph{Joint Genus and Degree Asymptotics for Connected Hurwitz Numbers}. No individual proposer is supplied; the proposed year is recorded as 2022.
\item Supplied report citation: \emph{Enumerative Combinatorics (2023)}, Oberwolfach Reports, pp.\ 3215--3305, DOI \texttt{10.4171/owr/2022/57}. Supplied source URL and literature-source link labelled ``original OWR report'': \url{https://doi.org/10.4171/owr/2022/57}. The report is the supplied authoritative reference; its contents have not been independently consulted for this draft.
\item The excerpt compares the requested asymptotics with work on triangulations of surfaces cited as [1], and attributes a signed-bijection proof for \(\ell=2\) to Fang in [3]. Full bibliographic entries for [1] and [3] are absent from the supplied record.
\item Literature assessment dated \texttt{2026-08-22}: the supplied curation labels the problem open and states that the simultaneous connected Hurwitz asymptotics and the proposed combinatorial understanding of the relevant coefficients remain unresolved in the identified literature. This is supplied dated curation, not an independently verified literature assessment. The background explicitly requests independent verification of time-sensitive status and cited claims.
\end{itemize}
\end{document}
